\scanpage{625}
\mspage{616}

% L1
Soit maintenant
\[
\tag{\unsure{20}}
\mu_{\tau'}=\{1,\tau'\}\subset\mathcal{G}^{\wedge}
\]
% L3
alors \inserted{\unsure{$\mu_{\tau'}$ normalise $L_\rho$}},
\deleted{\unsure{ce}} car
\(\tau'(\rho)=\rho^{-1}\) d'où \(\tau'(L_\rho)=L_\rho\)
donc il normalise \(Z_\rho\unsure{\equiv|=}\operatorname{Cent}_{\mathcal{M}}(L_\rho)\), donc
aussi \(Z_\rho^!=Z_\rho\cap\mathcal{M}^!\)
(puisque \(\mathcal{M}^!\) est invariant \unsure{dans} \(\mathcal{M}\))
\(=Z_{\rho,0}^!\times\mu=\mathcal{M}'(-\bar{\jmath})\times\mu\).
% L9
\deleted{Soit} D'ailleurs
\(\mu_{\tau'}\cap Z_\rho^!\subset
\deleted{\unsure{\{1\}}}\mu_{\tau'}\cap\mathcal{M}^!=\{1\}\).
Soit
\[
\tag{21}
\mathcal{N}_\rho^?=\mu_{\tau'}\cdot Z_\rho^!
=\mu_{\tau'}\cdot
\underbrace{(Z_{\rho,0}^!\times\mu)}_{
\substack{\mathcal{M}(-\bar{\jmath})\times\mu\\
\simeq\Gamma_{\mathbb{Q}}^{\prime\sim}\times\mu}}
\quad\deleted{\unsure{\simeq\mathcal{M}_{\mathbb{Q}}\times\mu|\simeq\mathcal{M}_{0}\times\mu}}.
\]
\hfill (produit $\frac12$ direct)

% L13
On fera attention que contrairement à ce
que pourrait suggérer cette écriture,
\(\mu_{\tau'}\) ne normalise \uline{pas} \(Z_{\rho,0}^!\) --- sinon on
aurait une \unsure{décomposition|disposition} en produit
\[
\mathcal{N}_\rho^?=\mathcal{N}_{\rho,0}^?\times\mu,
\quad\text{où }\mathcal{N}_{\rho,0}^?
=\deleted{\unsure{Z_{\rho,0}^!\times\mu}}
\mu_{\tau'}\cdot\unsure{Z_{\rho,0}^?|Z_{\rho,0}^!},
\]
% L19
or il n'en est rien. On a une suite exacte
\[
\tag{22}
\diagram{\begin{tikzcd}[ampersand replacement=\&,column sep=large,row sep=small]
1\arrow[r] \& \mu\arrow[r] \& \mathcal{N}_\rho^?\arrow[r]
\& \gamma_{\rho,\sigma}\arrow[r,no head]\arrow[d,equal] \& 1\\
\& \& \& \Gamma_{\mathbb{Q}}^{\sim} \&
\end{tikzcd}}
\]
% L22
de sorte que \(\mathcal{N}_\rho^?\) est une quasi-section
de multiplicité $2$ de l'extension
\(\mathcal{N}_\rho^{*}\longrightarrow\gamma_{\rho,\sigma}\), --- de
\(\mathcal{M}\longrightarrow\gamma_{\rho,\sigma}=\Gamma_{\mathbb{Q}}^{\sim}\).
Ainsi on trouve

\margin{624}

\scanpage{626}

% L1
\[
\tag{23}
\left\{\begin{aligned}
\mathcal{M}&=\mathcal{N}_\rho^?\wedge_{\mu_2}\mathcal{G}^{+\wedge}\\
\unsure{\mathcal{N}_\rho^*}&=\mathcal{N}_\rho^?\wedge_{\mu_2}
\deleted{\unsure{\mathbb{Z}}}L_\rho
\simeq\mathcal{N}_\rho^?\times
\underbrace{L_{-\rho}}_{\simeq\mathbb{Z}/3\mathbb{Z}}
\end{aligned}\right.
\]
% L3
La \deleted{suite}\inserted{extension} (22) de
\(\gamma_{\rho,\sigma}\) par \(\mu\) est
l'extension notée \(\widetilde{\gamma}_{\rho\sigma}\) dans la section précédente,
elle est sûrement \uline{\unsure{non} \deleted{triviale}\inserted{\unsure{scindée}}}.
Elle donne naissance à un \unsure{hom} canonique
\[
\tag{24}
\diagram{\begin{tikzcd}[ampersand replacement=\&,column sep=large]
\gamma_{\rho,\sigma}\arrow[r,"\sim"]
\&\mathcal{N}_\rho^?/\mu\arrow[r]
\&\mathcal{N}_\rho^*/\mu.
\end{tikzcd}}
\]

% L9
\uline{Remarque} Pour bien faire, il faudrait
expliciter l'opération de \(\tau'\) sur
\(Z_\rho^!\simeq Z_{\rho,0}^!\times\mu
\xrightarrow{\sim}\Gamma_{\mathbb{Q}}^{\prime\sim}\times\mu\), en termes de
l'opération de l'image \(\tau_{\Gamma\deleted{\unsure{\mathbb Q}}}\) de \(\tau'\) sur
\(\Gamma_{\mathbb Q}^{\prime\sim}\) par \unsure{autom.} int.
(NB \(\tau_\Gamma\) n'est autre que l'élément de
\(\Gamma_{\mathbb Q}\subset\Gamma_{\mathbb Q}^{\sim}\), ``conjugaison complexe'')
et d'un \unsure{hom} de \(\Gamma_{\mathbb Q}^{\prime\sim}\) dans \(\mu\)
\inserted{i.e d'un ``caractère quadratique'' de \(\Gamma_{\mathbb Q}^{\prime\sim}\)}
--- il y a fort à parier que
celui-ci est justement \(\varepsilon_2\) --- de sorte
que \(\Gamma_{\mathbb Q}^{0\sim}\simeq
\Gamma_{\mathbb Q}^{\prime\sim}\cap\Gamma_{\mathbb Q}^{\prime\prime\sim}\)
apparaît comme invariant par \(\tau'\) \ldots
En somme, il y aurait lieu de reprendre l'étude
de l'extension faite dans la section précédente,
en \unsure{y partant} du point de vue des \(\alpha,\beta\) \unsure{associés},

\margin{625}

\scanpage{627}
\mspage{617}

% L1
dans l'optique actuelle, sans doute plus
simple. Il y aurait lieu en $\overline{\mathrm m}$ temps
de \deleted{définir}\inserted{expliciter} une autre section partielle,
qui \deleted{mériterait} serait au dessus de
\(\deleted{\unsure{\Gamma}}\Gamma_{\mathbb Q}^{\prime\prime\prime\sim}\),
\[
\Gamma_{\mathbb Q}^{\prime\prime\prime\sim}
=\{x\in\Gamma_{\mathbb Q}^{\sim}\mid
\varepsilon_2(x)\varepsilon_3(x)=1\},
\]
% L6
et qui mériterait peut-être la notation
\(\mathcal{M}(\sqrt3)\) \unsure{ou du moins}
\(\mathcal{M}(q)\), pour
un générateur \(\unsure{q|\rho}\) convenable de \(\mathbb{Q}(\sqrt3)\) ???
\margin{?\quad\uline{\unsure{Méditer ces}}
\uline{\unsure{coïncidences.} \unsure{Attentif.|attentiv.}}}

% L9
Je me rends compte que je suis en train
de \unsure{déconner} --- ayant un peu perdu contact
avec le contenu géométrique de mes
calculs. Il faut garder à l'esprit que
\[
\tag{25}
\diagram{\begin{tikzcd}[ampersand replacement=\&]
\mathcal{M}_0^!\arrow[r,"\sim"] \& \mathcal{M}_{0,3}^{!\sim}
\end{tikzcd}}
\]
% L14
correspond \unsure{au} $\sim$ près, au \(\pi_1\) de \(U_{0,3\mathbb Q}\),
les sections de \(\mathcal{M}_0^!\longrightarrow
\Gamma_{\mathbb Q}^{\sim}\simeq\gamma_{\rho,\sigma}\) correspondent
donc (au $\sim$ près) à celles de
\(\pi_1(U_{0,3\mathbb Q})\) sur \(\Gamma_{\mathbb Q}\),
donc (``\unsure{invariablement|inévitablement|inavouablement}'') aux points
\inserted{\unsure{rat./$\mathbb Q$}} de \(U_{0,3\mathbb Q}\)
--- quand ce sont des sections partielles, aux
points de \(U_{0,3\mathbb Q}\) à valeurs dans des extensions
finies de \(\mathbb Q\). Ainsi, \deleted{les} la section
\deleted{\unsure{privilégiée}}\inserted{envisagée}
dans b) correspond bien au \(\mathrm{pt}\ \frac12\) de \(U_{0,3\mathbb Q}\),
la \deleted{\unsure{quasi-}}section
\inserted{partielle \(\mathcal M(-\bar\jmath)\)} envisagée ici dans c) correspond au
point \(-\bar\jmath\) de \(U_{0,3\mathbb Q}\), qui est dans
\(U_{0,3\deleted{\unsure{\mathbb Q}}}(\mathbb Q[\sqrt[3]{1}])\),
donc la notation \inserted{\(\mathcal M(-\bar\jmath)\)} est raisonnable.
Quant à \(\mathcal N_\rho^{\sim}\),
elle n'est \uline{pas} \(\subset\mathcal M_0^!\) --- $\overline{\mathrm m}$ en divisant par \(\mu\);
on trouve après division par \(\mu\) une section

\margin{626}

\scanpage{628}

% L1
de \(\mathcal M_{0,3}^{\sim}\) sur \(\Gamma_{\mathbb Q}^{\sim}\), mais ce n'est pas une
section de \(\mathcal M_{0,3}^{!\sim}\). Or au $\sim$ près, \(\mathcal M_{0,3}^{\sim}\)
est le \(\pi_1\) de \(M_{1,1\mathbb Q}'\), schéma modulaire sur \(\mathbb Q\)
des courbes elliptiques \uline{mod symétrie}.
Donc la donnée de cette section correspond à
la donnée d'une telle ``courbe elliptique
mod symétrie'' sur \(\mathbb Q\). Passant au \deleted{sous}
sous-groupe \(\Gamma_{\mathbb Q}^{\prime\sim}\) --- ce qui correspond à l'extension
du corps de base de \(\mathbb Q\) à \(\mathbb Q(j)=\mathbb Q(\sqrt[3]{1})\)
on retrouve l'image dans \(\mathcal M_{0,3}^{!\sim}\) de la \deleted{quasi-}
section \inserted{partielle} précédente --- \deleted{il} à \(\mathcal M(-\bar\jmath)\)
--- cela montre donc qu'il s'agit toujours de la ``$\overline{\mathrm m}$'' courbe
elliptique. J'ai donc l'impression \deleted{que} cependant
que sur \(\Gamma_{\mathbb Q}^{\sim}\) tout entier, cette section de
\(\mathcal M_{0,3}^{\sim}\) ne peut se remonter en une section
de \(\unsure{\mathcal M_\mu|\mathcal M}\) sur \(\Gamma_{\mathbb Q}^{\sim}\)
--- i.e justement que l'ext. \(\mathcal N_\rho^{\unsure{?|*}}\)
de \(\Gamma_{\mathbb Q}^{\sim}\) par \(\mu\) est non triviale.
Cela \uline{signifierait}-il qu'il n'existe pas de courbe elliptique
déf/\(\mathbb Q\) d'invariant $0$ ? (Chose que semble absurde
--- on doit pouvoir sur \uline{un corps} donner n'importe
quel invariant\ldots). \margin{?}
De façon précise, on a le
diagramme de schémas modulaires sur \(\mathbb Q\)
\[
\tag{26}
\diagram{\begin{tikzcd}[ampersand replacement=\&,column sep=large,row sep=large]
\begin{gathered}\text{sch. mod.}\\\text{des courbes}\\\text{ell.}\end{gathered}
\&M_{1,1}\arrow[d,"\substack{\text{gerbe}\\\text{de gr. }\mu}"]\arrow[dr]
\& \&\\
\begin{gathered}\text{sch. mod. des}\\\text{courbes ellipt.}\\\text{mod symétrie}\end{gathered}
\&M_{1,1}'\arrow[r]
\&\underset{\unsure{0|\circ}}{\mathbb E^1}
\&\begin{gathered}\text{schéma modulaire}\\\text{grossier}\end{gathered}
\end{tikzcd}},
\]
% L26
le fait qu'un point de \(M_{1,1}'(\mathbb Q)\) ne se re-

\margin{627}

\scanpage{629}
\mspage{618}

% L1
monte pas en un point de \(M_{1,1}\), ne \uline{signifie}
pas que son image dans \(\mathbb E^1\) ne se remonte pas
à \(M_{1,1}\). De façon précise, le \deleted{\illegible{1}}\inserted{diagramme}
s'identifie à :
\[
\tag{27}
\diagram{\begin{tikzcd}[ampersand replacement=\&,column sep=large,row sep=large]
M_{1,1}\simeq(U_{0,3},\widetilde{\mathcal G}_3)
  \arrow[d]\arrow[dr]\&\\
M_{1,1}'\simeq(U_{0,3},\mathcal G_3)\arrow[r]
  \&U_{0,3}/\mathcal G_3\simeq U_{0,3}/\widetilde{\mathcal G}_3
     \simeq\mathbb E^1
\end{tikzcd}},
\]
% L7
où \(\widetilde{\mathcal G}_3\) est l'extension de
\(\mathcal G_3\)\inserted{\(=\mathrm{SL}(2,\mathbb Z/2\mathbb Z)\)}
par \(\mu\), définie par
\[
\tag{28}
\begin{aligned}
\widetilde{\mathcal G}_3
 &=\mathcal G^+/\pi_0
 \simeq\mathrm{SL}(2,\mathbb Z)/
 \begin{gathered}
 \text{sous-groupe invariant}\\
 \text{engendré par}\\[-2pt]
 -\ell_0
 \end{gathered}
 \\
 &\simeq\mathrm{SL}(2,\mathbb Z/4\mathbb Z)/
 \begin{gathered}
 \text{sous-groupe}\\
 \text{invariant}\\
 \text{(d'ordre 4) engendré}\\
 \text{par }-\ell_0,\ \text{i.e \unsure{s\textsuperscript{s}groupe|s\textsuperscript{t}groupe}}\\
 \text{engendré par}\\
 -\ell_0,-\ell_1.
 \end{gathered}
\end{aligned}
\]
% L13
en faisant opérer \(\widetilde{\mathcal G}_3\) sur \(U_{0,3}\) par l'intermédiaire
de \(\mathcal G_3\). Un pt de \(M_{1,1}'\)\inserted{\unsure{rat/\(\mathbb Q\)}}
s'identifie \(\mathcal G_3\)-torseur \inserted{\(T\)} sur \(\mathbb Q\), muni
d'une application
\[
\tag{29}
\diagram{\begin{tikzcd}[ampersand replacement=\&]T\arrow[r]\&U_{0,3}\end{tikzcd}}
\]
\margin{\(T\) est le torseur
\inserted{\unsure{sous \(\mathrm{SL}(2,\mathbb Z/2\mathbb Z)\)}}
des \unsure{rigidifications|privilégiations} des points d'ordre $2$ et
\(\widetilde T\) celui des ``\unsure{pri-rigidifications|privilégiations|rigidifications}
des points d'ordre \(\unsure{\leq4|4}\)
\unsure{précisées}'' (le premier est défini
pour une courbe \unsure{ellipt.} mod symétrie, et
le $2^{\text{ième}}$ seulement pour une vraie courbe
elliptique)}
% L18
comp. avec les actions de \(\mathcal G_3\), un
point de \deleted{\(U_{0,3}\)} \(M_{1,1}\) s'identifie à un \(\widetilde{\mathcal G}_3\)-torseur
\(\widetilde T\) sur \(\mathbb Q\), muni d'une application
\[
\tag{30}
\diagram{\begin{tikzcd}[ampersand replacement=\&]\widetilde T\arrow[r,no head]\&U_{0,3}\end{tikzcd}}
\]
comp. avec l'opération de \(\widetilde{\mathcal G}_3\). \uline{Désignant}
par \deleted{\unsure{\(T\)}}
\[
T=\widetilde T/\mu
\]
le \(\mathcal G_3\)-torseur correspondant, la donnée de
(30) équivaut à la donnée de (29). Donc les

\margin{628}

\scanpage{630}

% L1
relèvements à \(M_{1,1}\) d'un point de \(M_{1,1}'\) rat/\(\mathbb Q\)
s'identifient aux \inserted{\(\widetilde{\mathcal G}_3\)-}torseurs \(\widetilde T\)
qui relèvent le \(\mathcal G_3\)-torseur \(T\), i.e. (à isom près) ils correspondent
aux hom \(\Gamma_{\mathbb Q}\longrightarrow\widetilde{\mathcal G}_3\)
qui relèvent l'hom
\(\Gamma_{\mathbb Q}\longrightarrow\mathcal G_3\) qui définit \(T\) :
\margin{on suppose \unsure{maintenant} choisi
un pt \deleted{\unsure{géom.}} de \(T(\overline{\mathbb Q}_0)\)}
\[
\diagram{\begin{tikzcd}[ampersand replacement=\&,column sep=large,row sep=small]
\&\widetilde{\mathcal G}_3\arrow[d]\\
\Gamma_{\mathbb Q}\arrow[r]\arrow[ur,dashed]\&\mathcal G_3
\end{tikzcd}}
\quad -
\]
% L8
l'obstruction est, comme il se doit, une
extension de \(\Gamma_{\mathbb Q}\) par \(\mu_2\), i.e. un \unsure{élém}
de \(H^2(\mathbb Q,\mu_2)\simeq {}_2\mathrm{Br}(\mathbb Q)\).
Dans le cas qui nous occupe, l'image de \(T\) dans \(U_{0,3}\)
est formé de \(\{-\bar\jmath,\unsure{-j|-\bar\jmath}\}\)
qui est connexe sur \(\mathbb Q\),
\uline{donc \(\Gamma_{\mathbb Q}\longrightarrow\mathcal G_3\) est surjectif}.
\margin{\unsure{corrigé|erroné}, cf rectification plus bas}
De façon générale, \deleted{\unsure{on doit}} dans la traduction
galoisienne en termes de scindages
de \(\mathcal M_{0,3}\) sur \(\Gamma_{\mathbb Q}\), l'hom
\(\Gamma_{\mathbb Q}\longrightarrow\mathcal G_3\)
est le composé
\[
\diagram{\begin{tikzcd}[ampersand replacement=\&]
\Gamma_{\mathbb Q}\arrow[r,"\unsure{k|\kappa}"]
\&\mathcal M_{0,3}\arrow[r]\&\mathcal G_3
\end{tikzcd}},
\]
% L22
qui est trivial ssi \deleted{\unsure{\(\Gamma_{\mathbb Q}\)}}
\inserted{\(\unsure{k|\kappa}\)} est une section de
\deleted{de}
\[
\mathcal M_{0,3}^!=\operatorname{Ker}
 (\mathcal M_{0,3}\longrightarrow\mathcal G_3)
\]
sur \(\Gamma_{\mathbb Q}\). Ici on a une section de
\(\mathcal M_{0,3}^!\) \deleted{sur}\inserted{au dessus}
du sous-groupe \(\Gamma_{\mathbb Q}'\) d'indice $2$,
\deleted{donc l'image de \(\Gamma_{\mathbb Q}'\)}
\(\unsure{k|\kappa}\mid\Gamma_{\mathbb Q}'=1\), donc
\(\unsure{k|\kappa}(\Gamma_{\mathbb Q})\subset\mathcal G_3\)
est d'ordre $2$, on voit que

\margin{629}

\scanpage{631}
\mspage{619}

c'est le sous-groupe \inserted{$\{1,\dot{\sigma}_{\infty}\}$} de $\mathcal{G}_3$.
(car l'image de $\tau'\inserted{{}=\sigma\tau}\in\mathcal{M}_{0,3}$ dans $\mathcal{G}_3$
\deleted{\unsure{est}} (comme celle de $\tau$ est $1$) est égale à celle de
$\sigma=\sigma_{\infty}$, i.e. $\dot{\sigma}_{\infty}$. Cela signifie
que \deleted{\unsure{T}} l'on a
\[
T\simeq T_0\mathbin{\wedge}_{\{1,\dot{\sigma}_{\infty}\}}\mathcal{G}_3
\]
i.e. $T$ est déduit d'un $\mu$-torseur (i.e. d'une
ext. quadratique de $\mathbb{Q}$\unsure{)} --- savoir justement
$\mathbb{Q}(\sqrt[3]{1})=\mathbb{Q}(j))$ par extension des groupes
d'opérateurs \inserted{$\mu\longrightarrow\mathcal{G}_3\quad(-1\longmapsto\dot{\sigma}_{\infty})$}.
Donc, prenant l'image inverse de $\mu=\{1,\dot{\sigma}_{\infty}\}$ dans
$\widetilde{\mathcal{G}}_3$, on trouve une extension de $\mu$ par $\mu$
--- ça ne peut guère être que $\mathbb{Z}/4\mathbb{Z}$ ! --- que j'ai envie de noter $\mathcal{E}''$

\[
\diagram{\begin{tikzcd}[ampersand replacement=\&,column sep=1.3em,row sep=0.15em]
\&\& \mathbb{Z}/4\mathbb{Z}\arrow[d,equal] \&\&\\
1\arrow[r]\&\mu\arrow[r]\&\mathcal{E}''\arrow[r]\&\mu\arrow[r,no head]\&1
\end{tikzcd}}
\]

\margin{oui, c'est évident puisque $\sigma=\sigma_{\infty}$ dans
$\widetilde{\mathcal{G}}_3=\mathrm{SL}(2,\mathbb{Z})/\unsure{\Pi_0|\Pi}$
est bien d'ordre $4$, donc $\mathcal{E}''\simeq\mathbb{Z}/4\mathbb{Z}$
\unsure{canoniquement}\ldots}
et le \unsure{pb} d'obstruction est de trouver un remontage

\[
\diagram{\begin{tikzcd}[ampersand replacement=\&,column sep=3em,row sep=2em]
\mathcal{E}''\arrow[r,no head]\&\mu\\
\Gamma_{\mathbb Q}\arrow[u,dashed]\arrow[ur]\&
\end{tikzcd}}
\]
,
(qui \deleted{\unsure{l'on}} n'existe sûrement pas) --- donc
l'obstruction s'identifie comme l'extension de $\Gamma_{\mathbb Q}$ par $\mu$,
image inverse de l'extension $\mathcal{E}''$.
\margin{C'est évident en \unsure{effet|fait}, \unsure{qu'on|puisque} :
$(\Gamma_{\mathbb Q})_{\mathrm{ab}}\simeq\widehat{\mathbb Z}^{*}
\supset\mathbb Z_3^{*}$, or
$(\mathbb Z_3^{*})_4\simeq\mathbb Z/2\mathbb Z$
(\unsure{donc déjà sur} $\mathbb Z_3^{*}$ \unsure{ça ne se remonte pas}\ldots}
(Il faudrait regarder pourquoi, en termes

\margin{630}

\scanpage{632}

classiques, cette obstruction est \unsure{scindée}, quand on
passe au sous-groupe $\Gamma_{\mathbb Q}^{\prime\prime\prime}$ correspondant au
corps $\mathbb Q(\sqrt{3}))$

Ainsi, le point $\mathbb Q$-rationnel envisagé de
$M_{1,1}'$ ne se remonte pas en un pt $\mathbb Q$-rationnel de
$M_{1,1}$ --- Par contre, son image $\inserted{0}$ dans
$\mathbb E^1\simeq U_{0,3}/\mathcal G_3$ doit bien se
remonter --- i.e. il doit bien exister un
$\widetilde{\mathcal G}_3$-torseur $\widetilde P$, et un
$\widetilde{\mathcal G}_3$-hom
\deleted{$\unsure{P}$}

\[
\diagram{\begin{tikzcd}[ampersand replacement=\&,column sep=1.6em,row sep=0.1em]
\widetilde P\arrow[r]\&U_{0,3}\&\text{i.e.}\&P\arrow[r]\arrow[d,equal]\&U_{0,3}\\
\&\&\&\widetilde P/\mu\&
\end{tikzcd}}
\]
.

Considérons en effet le \inserted{sous-}schéma
\deleted{de $U_{0,3}$, image inverse}
\[
\deleted{\unsure{Q=P}}\qquad
Q=U_{0,3}\mathbin{\times}_{\mathbb E^1}\{0\}
\]
l'image inverse de $0$, soit
\[
P_0=Q_{\mathrm{red}}\ \deleted{\unsure{\simeq}}
\]
c'est un sous-schéma qui, sur un corps
$K\supset\mathbb Q(j)$, \deleted{se scinde en} est somme de
deux pts, et qui sur $\mathbb Q$ est connexe
--- c'est un schéma \inserted{can.\textsuperscript{l}} isomorphe à
$\operatorname{Spec}\mathbb Q(-\bar j)$. Il est stable par $\mathcal G_3$,
qui y opère via

\[
\diagram{\begin{tikzcd}[ampersand replacement=\&]
\mathcal G_3\arrow[r,"\mathrm{sg}"]\&\mu
\end{tikzcd}}
\]

et l'opération galoisienne de $\mu$ sur le $\mu_{\mathbb Q}$-torseur $P_0$.

\margin{631}

\scanpage{633}
\mspage{620}

Ainsi on a le diagramme

\[
\diagram{\begin{tikzcd}[ampersand replacement=\&,column sep=1.8em,row sep=2em]
\&\widetilde{\mathcal G}_3\arrow[r,"\mathrm{can}"]\&\mathcal G_3\arrow[r,"\mathrm{sg}"]\&\mu\\
(31)\&\&\Gamma_{\mathbb Q}\arrow[ur]\&
\end{tikzcd}}
\]

\margin{\uline{NB} le noyau de
$\widetilde{\mathcal G}_3\longrightarrow\mu$ est une
extension \inserted{centrale} de $\mathbb Z/3\mathbb Z$
par $\mu=\mathbb Z/2\mathbb Z$, elle
est donc \unsure{trivial} et
isom. à $\mathbb Z/6\mathbb Z$
(noyau engendré par $\rho$).}

On voit donc que la catégorie des relèvements
du pt $\mathbb Q$-rationnel $0$ de $\mathbb E^1$ en un
pt $\mathbb Q$-rationnel de $M_{1,1}'$ resp. $M_{1,1}$, est
équivalente à celle des $\mathcal G_3$-torseurs \inserted{$P$} resp\unsure{.}
des $\widetilde{\mathcal G}_3$-torseurs \inserted{$\widetilde P$} sur $\mathbb Q$, qui relèvent le
$\mu$-torseur $P_0$. On avait choisi un
$P=T$ de façon \unsure{particulièrement} \unsure{trivial|triviale},
en utilisant le scindage

\[
\diagram{\begin{tikzcd}[ampersand replacement=\&,column sep=1.5em]
\mu\arrow[r,"\sim"]\&\{1,\dot\sigma_{\infty}\}\arrow[r,hook]\&\mathcal G_3
\end{tikzcd}}
\]

de $\mathcal G_3\longrightarrow\mu$ --- mais le torseur $T$ ne se
remonte pas en un $\widetilde T$. Les classes d'isomorphie de pts de $M_{1,1}'$ resp. de $M_{1,1}$ au-dessus
du pt $0$ de $\mathbb E^1$, \unsure{correspondent} de $\overline{\mathrm m}$ aux
classes de conjugaison \inserted{(mod noyaux \inserted{$\mathcal G_3^+,\widetilde{\mathcal G}_3^+$} de $\mathcal G_3\longrightarrow\mu$ et $\widetilde{\mathcal G}_3\longrightarrow\mu$)} de relèvements de
$\Gamma_{\mathbb Q}\longrightarrow\mu$ en $\Gamma_{\mathbb Q}\longrightarrow\mathcal G_3$ resp. en $\Gamma_{\mathbb Q}\longrightarrow\widetilde{\mathcal G}_3$.

\deleted{\unsure{De là}} Considérons un relèvement en

\[
\diagram{\begin{tikzcd}[ampersand replacement=\&]
\Gamma_{\mathbb Q}\arrow[r]\&\mathcal G_3
\end{tikzcd}}
\]
,
son image est un sous-groupe de $\mathcal G_3$, qui est soit

\margin{632}

\scanpage{634}

$\{1,\dot\sigma_{\infty}\}$ ou un de ses conjugués,
$\{1,\dot\sigma_0\}$, $\{1,\unsure{\bar\sigma_1|\dot\sigma_1}\}$,
\uline{ou} c'est $\mathcal G_3$ tout entier. Dans le premier
cas, on a vu que ça ne se remonte pas
en $\widetilde{\mathcal G}_3$. Donc pour pouvoir remonter,
il faut prendre un \uline{épi} $\Gamma_{\mathbb Q}\longrightarrow\mathcal G_3$ ---
on sort donc du cadre abélien sur
$\mathbb Q$, pour entrer dans le domaine des
extensions abéliennes de $\mathbb Q(j)$.

Examiner en termes des groupes de
Galois s'il existe un tel relèvement \deleted{qui}
$\Gamma_{\mathbb Q}\longrightarrow\widetilde{\mathcal G}_3$, et \deleted{\unsure{quel}}
comment les obtenir,
apparaît comme un exercice plaisant
et \unsure{délectable} de théorie des corps de
classes, que je ne vais pas poursuivre
pour le moment, pour ne pas éterniser
cette digression. Revenant par la suite
sur ce point, il me faudrait en $\overline{\mathrm m}$
temps expliciter la relation entre les
\unsure{prises|points} de un remontage d'l'hom de $\Gamma_{\mathbb Q}$,
et scindages de l'extension $\mathcal M_{0,3}$ de $\Gamma_{\mathbb Q}$
par $\unsure{\mathcal G_{0,3}^{+\wedge}|\widetilde{\mathcal G}_{0,3}^{+\wedge}}$,
ou $\mathcal M$ de $\Gamma_{\mathbb Q}$ par
$\mathcal G^{+\wedge}\simeq\mathrm{SL}(2,\mathbb Z)^\wedge$.

\margin{633}

\scanpage{635}
\mspage{621}

Revenons aux sous-groupes \deleted{$\unsure{\mathcal N_{\rho}^{?}}$ de}

\[
\diagram{
\begin{tikzcd}[ampersand replacement=\&,column sep=3em,row sep=2em]
\mathcal N_{\rho}^{?}
  \arrow[r,phantom,"{\overset{3}{\unsure{\subset}}}"{description}]
\&\mathcal N_{\rho}^{*}\arrow[r,phantom,"{\subset}"{description}]
\&\mathcal M\\
Z_{\rho,0}^{!}=\mathcal M'(-\bar j)
  \arrow[u,phantom,"{4\,\cup}"{description}]
\&\&
\end{tikzcd}
}
\]

contenant \unsure{la|le} \deleted{section}\inserted{\unsure{sous-groupe|sous-groupes}}
partielle $\mathcal M(-\bar j)\simeq\widetilde\Gamma_{\mathbb Q}'$.

Soit
\[
\tag{32}
\begin{aligned}
\mathcal M^0(-\bar j)&\overset{2}{\subset}\mathcal M'(-\bar j)\\
\underset{\unsure{\mathrm{dfn}|\mathrm{def}}}{=}&\operatorname{Ker}
\bigl(\mathcal M'(-\bar j)\xrightarrow{\varepsilon_2}\mu\bigr)
=\deleted{\unsure{\mathcal N_{\rho}^{?}|\mathcal M}}\,Z_{\rho}
\cap\deleted{\unsure{\operatorname{Ker}\varepsilon_2}}\,
\mathcal M_0^{\prime\prime!}\\
&=\mathcal N_{\rho}^{*}\cap\mathcal M_0^{0!}
\end{aligned}
\]
qui est donc une section partielle \unsure{en|au} dessus de
\[
\tag{33}
\begin{aligned}
\widetilde\Gamma_{\mathbb Q}^0
&=\operatorname{Ker}\bigl(\widetilde\Gamma_{\mathbb Q}
\xrightarrow{(\varepsilon_2,\varepsilon_3)}\mu\times\mu\bigr)
\deleted{\unsure{\ne|=}\widetilde\Gamma_{\mathbb Q}}\\
&=\operatorname{Ker}\bigl(\widetilde\Gamma_{\mathbb Q}
\longrightarrow\widehat{\mathbb Z}^{*}\longrightarrow
(\mathbb Z/12\mathbb Z)^{*}\bigr)
\end{aligned}
\]

\deleted{Considérons} Modulo une vérification
que je vais faire plus bas, $\tau'$ normalise
\deleted{(34)}
$\mathcal M^0(-\bar j)$, considérons \unsure{les|le}
\[
\tag{34}
\mathcal M'''(-\bar j)\deleted{\unsure{\mathcal N_{\rho,0}^{?}}}
\overset{\unsure{\mathrm{dfn}|\mathrm{def}}}{=}
\underset{\text{(semi-direct)}}{\mu_{\tau'}\mathbin{.}\mathcal M^0(-\bar j)}
\qquad\text{où }\mu_{\tau'}=\{1,\tau'\}
\]
c'est un sous-groupe d'indice $4$ de $\mathcal N_{\rho}^{?}$,
et on a maintenant
\[
\tag{35}
\deleted{\unsure{\mathcal N_{\rho,0}^{?}}}\quad
\mathcal M'''(-\bar j)\xrightarrow{\sim}\widetilde\Gamma_{\mathbb Q}''',
\]
c'est une section partielle \unsure{en|au} dessus
de $\widetilde\Gamma_{\mathbb Q}'''$, d'où des isomorphismes

\margin{634}

\scanpage{636}

\deleted{
\[
\tag{36}
\left\{
\begin{aligned}
\mathcal M'''&\simeq\mathcal M'''(-\bar j)\mathbin{.}\mathcal G^{+\wedge}\\
\deleted{\unsure{\mathcal M'''}}\,\mathcal N_{\rho}^{*}
&\simeq\mathcal M'''(-\bar j)\mathbin{.}L_{\rho}.
\end{aligned}
\right.
\]

On a le diagramme récapitulatif
d'inclusions de sous-groupes de $\mathcal N_{\rho}^{*}$

\[
\diagram{
\begin{tikzcd}[ampersand replacement=\&,column sep=2.5em,row sep=1.5em]
\&\&\&Z_{\rho}\arrow[dr,hook,"2"]\&\\
\&\&\begin{gathered}\mathcal M'(-\bar j)\\=Z_{\rho,0}^{!}\end{gathered}
  \arrow[ur,hook,"6"]\arrow[dr,hook,"4"]
\&\&\mathcal N_{\rho}^{*}\\
(37)\&\mathcal M^0(-\bar j)\arrow[ur,hook,"2"]\arrow[dr,hook,"2"']
\&\&\mathcal N_{\rho}^{?}\arrow[ur,hook,"3"']\&\\
\&\&\mathcal M'''(-\bar j)\arrow[ur,hook,"4"]\&\&
\end{tikzcd}
}
\]


Notons que l'inclusion
\[
\tag{38}
\underset{\substack{\simeq\\\unsure{\mathbb D_6|\mathbb D_{\sigma}}}}{\mathbb D_{\rho}}
\overset{\mathrm{dfn}}{=}
\underset{\simeq\mathbb Z/2\mathbb Z}{\uline{\mu_{\tau'}}}
\mathbin{.}
\underset{=\mathbb Z/6\mathbb Z}{\uline{L_{\rho}}}
\subset\mathcal N_{\rho}^{*}
\]
induit un isomorphisme d'
\[
\tag{39}
\mathbb D_{\rho}\xrightarrow{\sim}\mathcal N_{\rho}^{*}/
\]
}

\margin{635}

\scanpage{637}
\mspage{622}

\[
\tag{36}
\left\{
\begin{aligned}
\mathcal{M}'''&\simeq \mathcal{M}'''(-\bar j)\mathbin{.}\mathcal{G}^{+\wedge}\\
\mathcal{N}_{\rho}^{*}&\simeq \mathcal{M}'''(-\bar j)\mathbin{.}L_{\rho}.
\end{aligned}
\right.
\]

Diagramme récapitulatif des sous-groupes de $\mathcal{N}_{\rho}^{*}$
obtenus jusqu'à présent:


\[
\diagram{
\begin{tikzcd}[ampersand replacement=\&,column sep=1.05em,row sep=2.8em]
\&\& \mathcal{M}^{0}(-\bar j)
  \arrow[rr,hook,"2"]
  \arrow[dl,hook,"2"']
\&\& \begin{gathered}\mathcal{M}'(-\bar j)\\=Z_{\rho,0}^{!}\end{gathered}
  \arrow[rr,hook,"2"]
  \arrow[dl,hook,"2"']
  \arrow[dr,no head,draw=gray,
    "{\unsure{\mathcal{M}'(-\bar j)\times\mu}}"{description,pos=.5},
    "{\unsure{\searrow}}"{pos=.93}]
\&\& \mathcal{M}'(-\bar j)\times\mu
  \arrow[rr,hook,"3"]
  \arrow[dl,hook,"2"']
\&\& Z_{\rho}\arrow[dl,hook,"2"']\\
(37)\&\mathcal{M}'''(-\bar j)\arrow[rr,hook,"2"]
\&\& \mathcal{M}'''(-\bar j)\times\mu\arrow[rr,hook,"2"]
\&\& \mathcal{N}_{\rho}^{?}\arrow[rr,hook,"3"]
\&\& \mathcal{N}_{\rho}^{*}
\&\&\\
\&\&\&\&1
  \arrow[uu,hook]
  \arrow[uull,hook,start anchor=north west,end anchor=south west]
  \arrow[rr,hook,"2"]
  \arrow[dl,hook,"2"']
\&\&\mu\arrow[uu,hook]\arrow[rr,hook,"3"]\arrow[dl,hook,"2"']
\&\&L_{\rho}\arrow[uu,hook]\arrow[dl,hook,"2"']\\
\&\&\&\mu_{\tau'}\arrow[uu,hook]\arrow[uull,hook]\arrow[rr,hook,"2"]
\&\&\mathbb{D}_{\tau'}\arrow[uu,hook]\arrow[rr,hook,"3"]
\&\&\mathbb{D}_{\rho}\arrow[uu,hook]
\end{tikzcd}
}
\]
,

où nous sommes ici intéressés \unsure{surtout} à la
“face supérieure”, \unsure{contenant|concernant} des sous-groupes d'indice fini de
$\mathcal{N}_{\rho}^{*}$, intermédiaires entre celui-ci et
$\unsure{\mathcal{M}_{0}(-\bar j)|\mathcal{M}^{0}(-\bar j)}$
(d'indice $24$ dans $\mathcal{N}_{\rho}^{\unsure{*}}$).
Je vais examiner la question d'invariance de ces sous-groupes
les uns dans les Autres, et la structure des groupes quotients.

Je vais refaire un diagramme qui montre \unsure{vraiment|mieux} les
\deleted{\unsure{groupes|rapports}} relations entre les groupes en question.

\margin{636}

\scanpage{638}


\[
\diagram{
\begin{tikzcd}[ampersand replacement=\&,column sep=.6em,row sep=1.3em,cells={font=\scriptsize}]
\&\&Z_{\rho}\arrow[rrr,hook,"2"]
  \arrow[ddddll,dashed,leftrightarrow,bend right=35,
    "{\unsure{\leftrightarrow}}"'{description}]
\&\&\&\mathcal{N}_{\rho}^{*}\\
\&\&\&Z_{\rho,0}^{!}=\mathcal{M}'(-\bar j)
  \arrow[rrr,hook,"2"]
  \arrow[dl,hook,"2"']
\&\&\&\mathcal{M}'(-\bar j)\mathbin{.}\mu_{\tau'}\;{}^{*}
  \arrow[dl,hook,"2"']\\
\&\&Z_{\rho}^{!}=\mathcal{M}'(-\bar j)\times\mu
  \arrow[uu,hook,"3"']
  \arrow[rrr,hook,"2"]
  \arrow[ddll,hook,dashed,"3"']
\&\&\&\mathcal{N}_{\rho}^{?}
  \arrow[uu,hook,"3"']
  \arrow[ddll,hook,dashed,"3"']\\
\&\&\&\mathcal{M}^{0}(-\bar j)
  \arrow[uu,hook,"2"']
  \arrow[rrr,hook,"2"]
  \arrow[ddl,hook,"2"']
\&\&\&\mathcal{M}'''(-\bar j)
  \arrow[uu,hook,"2"']
  \arrow[ddl,hook,"2"']\\
Z_{\rho}\arrow[rrr,hook,dashed]
\&\&\&\unsure{\mathcal{N}_{\rho}|\mathcal{N}_{\rho}^{*}}
  \arrow[uuuurr,dashed,leftrightarrow,bend left=20,
    "{\unsure{\leftrightarrow}}"'{description}]\\
\&\&\mathcal{M}^{0}(-\bar j)\times\mu
  \arrow[uuu,hook,"2"']
  \arrow[rrr,hook,"2"]
  \arrow[dddll,hook,
    "{\deleted{\unsure{\hookrightarrow}}}"{description}]
\&\&\&\mathcal{M}'''(-\bar j)\times\mu
  \arrow[uuu,hook,"2"']
  \arrow[dddll,hook,
    "{\deleted{\unsure{\hookrightarrow}}}"{description}]\\
\&\&\&1\arrow[uuu,hook]
  \arrow[rrr,hook,"2"]
  \arrow[dl,hook,"2"']
\&\&\&\mu_{\tau'}\arrow[uuu,hook]\arrow[dl,hook,"2"']\\
(38)\&\&\mu\arrow[uu,hook]\arrow[rrr,hook,"2"]
  \arrow[ddll,hook,dashed,"3"']
\&\&\&\mathbb{D}_{\tau'}\arrow[uu,hook]
  \arrow[ddll,hook,dashed,"3"']\\
\deleted{\unsure{L_{\rho}}}
  \arrow[rrr,hook,
    "{\deleted{\unsure{\hookrightarrow}}}"{description}]
\&\&\&\deleted{\unsure{\mathbb{D}_{\rho}}}\\
L_{\rho}\arrow[uuuuu,hook,dashed]
  \arrow[rrr,hook,dashed,"2"]
\&\&\&\mathbb{D}_{\rho}\arrow[uuuuu,hook,dashed]
\end{tikzcd}
}
\]


\margin{$*$ -- Attention,
$\mathcal{M}'(-\bar j)\mathbin{.}\mu_{\tau'}$ n'est \uline{pas}
un sous-groupe, car $\mu_{\tau'}$ ne normalise pas
$\mathcal{M}'(-\bar j)$.}

\unsure{a)} Sous-groupes d'indice fini dans $Z_{\rho}$ -- dans le diagramme
il y en a quatre


\[
\diagram{
\begin{tikzcd}[ampersand replacement=\&,column sep=3em,row sep=3em]
\&Z_{\rho}^{!}=\mathcal{M}'(-\bar j)\times\mu
\&\mathcal{M}'(-\bar j)=Z_{\rho,0}^{!}=\mathcal{N}_{\rho,0}^{\unsure{*!}}
  \arrow[l,hook,"2"']\\
(39)\&\mathcal{M}^{0}(-\bar j)\times\mu
  \arrow[u,hook,"2"']
\&\mathcal{M}^{0}(-\bar j)=Z_{\rho,0}^{0!}
  \arrow[u,hook,"2"]
  \arrow[l,hook,"2"]
\end{tikzcd}
}
\]


sont ils invariants dans $Z_{\rho}$? C'est clair pour les sous-groupes
\deleted{$\unsure{\mathcal{M}'(-\bar j)}$}
\[
\tag{40}
\left\{
\begin{aligned}
\mathcal{M}'(-\bar j)\times\mu
  &=Z_{\rho}\cap\mathcal{M}^{!}=Z_{\rho}^{!},
  \quad\text{et pour}\\
\deleted{\unsure{\mathcal{M}^{0}(-\bar j)\times\mu=}}&\\[-.5ex]
\mathcal{M}^{0}(-\bar j)\times\mu
  &=Z_{\rho}\cap\mathcal{M}^{!}\cap
       \operatorname{Ker}(\varepsilon_{2}:\mathcal{M}\to\mu)
\end{aligned}
\right.
\]
-- cela montre que ces groupes sont

\margin{637}

\scanpage{639}
\mspage{623}

même invariants dans $\mathcal{N}_{\rho}^{*}$, puisque
\deleted{$\unsure{\mathcal{M}^{!}}$}.
$Z_{\rho},\mathcal{M}^{!},\operatorname{Ker}\varepsilon_{2}$
sont stabilisés par $\mathcal{N}_{\rho}^{*}$.
D'ailleurs, \deleted{comme}
$Z_{\rho,0}^{!}=\mathcal{M}'(-\bar j)$ est
\deleted{d'indice $2$}\inserted{invariant}
\deleted{\unsure{dans}}\inserted{dans}
$Z_{\rho}^{!}=\mathcal{M}'(-\bar j)\times\mu$,
\unsure{pour} montrer qu'il l'est dans $Z_{\rho}$, il suffit de prouver
qu'il est normalisé par $\rho$
(puisque $Z_{\rho}=L_{\rho}\mathbin{.}Z_{\rho}^{!}$ -- c'est
$\pm$ trivial), or examinons de façon générale pour
\deleted{$\unsure{x\in}$} $x\in\mathcal{M}$
si $x$ normalise \deleted{$\unsure{\mathcal{M}_{0}^{!}}$}
\[
\mathcal{M}_{0}^{!}\overset{\mathrm{dfn}}{=}
\mathcal{M}(0)\mathbin{.}\Pi_{0},
\]
Comme il normalise $\Pi_{0}$, il suffit de voir si
pour $u\in\mathcal{M}(0)$, on a
\[
xux^{-1}\in\mathcal{M}_{0}^{!}
\]
ce qui équivaut aussi à
\[
xux^{-1}\equiv u\qquad(\Pi_{0})
\]
i.e.
\[
[x,u]\in\Pi_{0}
\]
or on sait que pour $x\in\mathcal{G}$,
\inserted{$u\in\mathcal{M}^{!}$,} on a
\[
\tag{*}
\underbrace{[x,u]}_{\inserted{x\,u(x)^{-1}}}^{\unsure{c}}
\equiv
(\operatorname{sg}(x))^{\inserted{\unsure{\xi_{2}(u)|\varepsilon_{2}(u)}}}
\deleted{\unsure{\xi_{2}(u)|\varepsilon_{2}(u)}}
\quad\mathrm{mod}\ \Pi_{0}
\]
où

\[
\diagram{
\begin{tikzcd}[ampersand replacement=\&,column sep=2em]
\operatorname{sg}:\&
\mathcal{G}\arrow[r,"\mathrm{can}"]\&
\mathcal{G}_{3}\arrow[r,"\operatorname{sg}"]\&
\mu.
\end{tikzcd}
}
\]


Cela montre \deleted{\unsure{qu'on}} que

\uline{Lemme}
\[
\left\{
\begin{aligned}
&\text{Si }x\in\mathcal{G}^{+},\ u\in\mathcal{M}_{0}^{!},
  \quad\text{alors }
  \deleted{xux^{-1}\in\mathcal{M}_{0}^{!}}
  \ \text{conditions équivalentes}\\
&a)\quad xux^{-1}\in\mathcal{M}_{0}^{!}\\
&b)\quad xux^{-1}\equiv u\quad(\Pi_{0})
       \quad\text{i.e.}\quad [x,u]\in\Pi_{0}\\
&c)\quad .\operatorname{sg}(x)^{
       \inserted{\unsure{\xi_{2}(u)|\varepsilon_{2}(u)}}}
       \deleted{\unsure{\xi_{2}(u)|\varepsilon_{2}(u)}}=1
       \quad\text{i.e.}\quad
       \operatorname{sg}(x)=1\quad\text{\uline{ou}}\quad
       \varepsilon_{2}(u)=1
\end{aligned}
\right.
\]

\margin{638}

\scanpage{640}

\deleted{Pour $x=\rho$, cela montre que \unsure{le|la}
\deleted{\unsure{condition}}\inserted{\unsure{premier groupe}}
\unsure{est}}

\uline{Corollaire 1} Soit $x\in\mathcal{G}^{+\wedge}$.
Pour que $x$ normalise $\mathcal{M}_{0}^{!}$,
il faut s. qu'on ait $\operatorname{sg}(x)=+1$.

\uline{Corollaire 2}. Le sous-groupe
$\mathcal{M}_{0}^{\prime\prime!}
=\mathcal{M}''(0)\mathbin{.}\Pi_{0}$
est \deleted{invariant dans} normalisé par
$\mathcal{G}^{+\wedge}$.

Ceci montre en particulier que $\rho$ normalise
$\mathcal{M}_{0}^{!}$, donc aussi
$\inserted{\mathcal{M}'(-\bar j)=}
Z_{\rho,0}^{!}=Z_{\rho}\cap\mathcal{M}_{0}^{!}$,
qui est donc invariant dans $Z_{\rho}$.
Il en est donc de $\overline{\mathrm m}$ de son intersection
avec $\mathcal{M}^{0}(-\bar j)\times\mu$, soit
$\mathcal{M}^{0}(-\bar j)$.
(Ou encore; ce groupe est égal à
$Z_{\rho}\cap\mathcal{M}_{0}^{\prime\prime!}$,
et $\mathcal{M}_{0}^{\prime\prime!}$ est normalisé par $\rho\ldots$).
\uline{Ainsi les quatre groupes} $(39)$
\uline{sont distingués dans $Z_{\rho}$.} )

Pour voir s'ils sont distingués dans
$\unsure{\mathcal{N}_{\rho}^{*}}
=\mu_{\tau'}\mathbin{.}Z_{\rho}$,
il suffit de voir s'ils sont normalisés par $\tau'$.
C'est clair pour $Z_{\rho}^{!}=Z_{\rho}\cap\mathcal{M}^{!}$
-- d'où le fait que
$\mathcal{N}_{\rho}^{?}=\mu_{\tau'}\mathbin{.}Z_{\rho}^{!}$
est un sous-groupe de $\mathcal{N}_{\rho}^{*}$.
Pour
\deleted{$\mathcal{M}^{0}(-\bar j)=
Z_{\rho}\cap\mathcal{M}_{0}^{\prime\prime!}$, il}
\inserted{voir \deleted{\unsure{si les}} les quels des}
trois autres sous-groupes
\inserted{envisagés}
en dernier sont normalisés, il faudrait donc expliciter
l'action de $\tau'$ \deleted{dessus.}
sur $Z_{\rho}^{!}=\mathcal{M}'(-\bar j)\times\mu$.

\margin{639}

\scanpage{641}
\mspage{623}

Notons que $\tau\in\mathcal{M}(0)\subset\mathcal{M}_{0}^{!}$
normalise $\mathcal{M}_{0}^{!}$, donc
\deleted{\unsure{on|donc}}\inserted{il} opère
\deleted{$\unsure{Z}$} sur
$\mathcal{M}^{!}
=\underbrace{\mathcal{M}_{0}^{!}}_{\simeq\widetilde{\Gamma}_{\mathbb{Q}}\mathbin{.}\Pi_{0}}
\times\mu$
par l'opération produit de la conjugaison par $\tau$ dans
$\mathcal{M}_{0}^{!}$, et l'identité dans $\mu$.
On a d'autre part $\tau'=\tau\sigma$, on a vu que
\[
 (40)\qquad
 \sigma(x)\equiv\varepsilon_{2}(x)x\qquad(\Pi_{0})
 \qquad x\in\mathcal{M}
\]
donc
\[
 (41)\qquad
 \tau'(x)\equiv\varepsilon_{2}(x)\tau(x)\qquad(\Pi_{0})
 \qquad x\in\mathcal{M}
\]
donc il s'ensuit
\[
 (42)\quad
 \left\{
 \begin{aligned}
 &\text{si }x\in\mathcal{M}_{0}^{!},\quad
   \text{alors }\tau'(x)\equiv x\quad(\Pi_{0})
   \inserted{\Longleftrightarrow\tau'(x)\in\mathcal{M}_{0}^{!}}
   \Longleftrightarrow\varepsilon_{2}(x)=1\\
 &\text{en particulier, }\mathcal{M}_{0}^{\prime\prime!}
   \text{ est normalisé par }\tau
 \end{aligned}
 \right.
\]
Donc \deleted{si $x\in$} appliquant ceci aux éléments de
$Z_{\rho,0}^{!}=\mathcal{M}'(-\bar j)$ dans
$Z_{\rho}^{!}=\mathcal{M}'(-\bar j)\times\mu$,
on trouve
\[
 (43)\quad
 \left\{
 \begin{aligned}
 &\text{Soit }x\in Z_{\rho,0}^{!}=\mathcal{M}'(-\bar j),
   \text{ alors }\deleted{\unsure{\tau}}\\
 &\tau'(x)\in Z_{\rho,0}^{!}
   \Longleftrightarrow\varepsilon_{2}(x)=1
 \end{aligned}
 \right.
\]

\uline{Corollaire}\quad
$Z_{\rho,0}^{!}\inserted{=\mathcal{M}'(-\bar j)}$
n'est pas normalisé par $\tau'$, donc n'est pas invariant dans
$\unsure{\mathcal{N}_{\rho}^{*}}$,
mais \deleted{le}
$\mathcal{M}^{0}(-\bar j)=Z_{\rho}\cap\mathcal{M}_{0}^{\prime\prime!}$
est normalisé par $\tau'$, donc aussi son produit par $\mu$.

Donc le seul des quatre sous-groupes $(39)$ de $Z_{\rho}$
qui ne soit distingué dans $\mathcal{N}_{\rho}$ est
$\mathcal{M}'(-\bar j)=Z_{\rho,0}^{!}$.
Les autres donnent

\margin{640}

\scanpage{642}

naissance, par produit $\frac12$ direct avec $\mu_{\tau'}$,
à des \unsure{s-groupes|ss-groupes} de $\mathcal{N}_{\rho}$,
invariants dans $\mathcal{N}_{\rho}$.
Le plus petit de tous ces sous-groupes est
$\mathcal{M}^{0}(-\bar j)$, qui est d'indice $24$ dans
$\unsure{\mathcal{N}_{\rho}^{*}}$.
Le groupe quotient \deleted{\unsure{est un}}
\[
 \deleted{\unsure{\text{groupe}}}\qquad
 \mathcal{N}_{\rho}^{*}/\mathcal{M}^{0}(-\bar j)
\]
est donc un groupe d'ordre $24$, qui reçoit d'ailleurs le groupe
$\mathbb{D}_{\rho}
\inserted{=\mu_{\tau'}\mathbin{.}L_{\rho}}
\subset\mathcal{N}_{\rho}^{*}$,
normalisateur de $L_{\rho}$ dans
$\unsure{\operatorname{GL}}(2,\mathbb{Z})$,
groupe isomorphe à $D_{6}$, et d'ordre $12$.
Le noyau de

\[
\diagram{
 \begin{tikzcd}[ampersand replacement=\&]
 (44)\quad\mathbb{D}_{\rho}
 \arrow[r,hook] \&
 \mathcal{N}_{\rho}^{*}/
 \mathcal{M}^{0}(\bar j)
 \end{tikzcd}
}
\]

est évidemment contenu dans $L_{\rho}=\mathbb{D}_{\rho}\cap Z_{\rho}$
et il est réduit à $1$ car
puisqu'on a $\overline{\mathrm m}$
\[
 L_{\rho}\cap Z_{\rho,0}^{!}
 \subset L_{\rho}\cap\mathcal{M}_{0}^{!}
 =L_{\rho}\cap\Pi_{0}=\{1\}.
\]
Ainsi $\mathbb{D}_{\rho}\inserted{\text{(d'ordre }12\text{)}}$
apparaît comme un sous-groupe d'indice $2$ de
$\mathcal{N}_{\rho}^{*}/\mathcal{M}^{0}(-\bar j)$,
tout comme il est sous-groupe d'indice $2$ de
$\mathcal{G}/\Pi_{0}$. Il serait tentant

\margin{641}

\scanpage{643}
\mspage{624}

alors d'établir un isomorphisme
\[
 \mathcal{N}_{\rho}^{*}/\mathcal{M}^{0}(-\bar j)
 \overset{?}{\simeq}\mathcal{G}^{\wedge}/\Pi_{0}
\]
compatible avec ces plongements de $\mathbb{D}_{\rho}$.
Essayons-le \unsure{ainsi|aussi}:

\uline{Proposition}\quad Considérons l'hom. composé

\[
\diagram{
 \begin{tikzcd}[ampersand replacement=\&,column sep=large,row sep=large]
 (45)\quad\mathcal{M}\arrow[r,"\Theta_{\mathcal{M}}"]\&
 \operatorname{GL}(2,\widehat{\mathbb{Z}})
 \arrow[r]\&
 \operatorname{GL}(2,\mathbb{Z}/4\mathbb{Z})
 \arrow[r,"\substack{\text{épi. de}\\\text{degré }4}"]\&
 \begin{gathered}
 \mathcal{G}^{\unsure{*|\wedge}}/\Pi_{0}\\
 \simeq\mathcal{G}^{\wedge}/\Pi_{0}^{\wedge}
 \end{gathered}\\
 \&\&\underbrace{\operatorname{GL}(2,\mathbb{Z})}_{\mathcal{G}}
 /\mathcal{G}[4]\arrow[u,"\sim"']\&
 \end{tikzcd}
}
\]

\deleted{et l'homomorphisme induit}
qui induit sur $\mathcal{G}^{\wedge}\subset\mathcal{M}$
l'homomorphisme canonique
$\mathcal{G}^{\wedge}\longrightarrow
\mathcal{G}^{\wedge}/\Pi_{0}^{\wedge}$.
Alors l'hom. \deleted{composé}\inserted{induit}

\[
\diagram{
 \begin{tikzcd}[ampersand replacement=\&]
 \mathcal{N}_{\rho}^{*}\arrow[r]\&\mathcal{G}/\Pi_{0}
 \end{tikzcd}
}
\]

a comme noyau $\mathcal{M}^{0}(-\bar j)$, et induit par passage
au quotient un isomorphisme

\[
\diagram{
 \begin{tikzcd}[ampersand replacement=\&]
 (46)\quad\mathcal{N}_{\rho}^{*}/\mathcal{M}^{0}(-\bar j)
 \arrow[r,"\sim"]\&\mathcal{G}/\Pi_{0}
 \end{tikzcd}
}
\]


Démonstration. Les éléments de $\mathcal{M}^{0}(-\bar j)$
sont les éléments de la forme
\[
 \left\{
 v=\alpha^{-1}u\ \middle|\
 \begin{aligned}
 &u\in\mathcal{M}(0),\quad\alpha\in\Pi_{0}\\
 &u(\rho)=\alpha(\rho)
   \qquad(\text{donc }\varepsilon_{3}(u)=1)\\
 &\varepsilon_{2}(u)=1
 \end{aligned}
 \right\}
\]
ils sont en corr. $1$-$1$ avec les éléments de
\[
 \mathcal{M}^{0}(0)
 =\{u\in\mathcal{M}(0)\mid\varepsilon_{2}(u)=\varepsilon_{3}(u)=1\}.
\]
L'image \unsure{de|des}

\margin{642}
